\chapter{2011年天津卷（文）}

\section{单选题}

\begin{problem}
\(\i\) 是虚数单位，复数 \(\frac{1 - 3\i}{1 - \i} =\)（\quad）

\choices
    {\(2 - \i\)}
    {\(2 + \i\)}
    {\(-1 - 2\i\)}
    {\(-1 + 2\i\)}
\end{problem}

\begin{answer}
A
\end{answer}

\begin{solution}
分子、分母同乘 \(1+\i\)，得
\[
\frac{1-3\i}{1-\i}
=\frac{(1-3\i)(1+\i)}{(1-\i)(1+\i)}
=\frac{4-2\i}{2}
=2-\i.
\]
故选 A.
\end{solution}

\begin{problem}
设变量 \(x\)，\(y\) 满足的约束条件 \(\begin{cases}
x\geqslant 1,\\
x+y-4\leqslant 0,\\
x-3y+4\leqslant 0
\end{cases}\)，则目标函数 \(z = 3x - y\) 的最大值为（\quad）

\choices
    {\(-4\)}
    {\(0\)}
    {\( \frac{4}{3} \)}
    {\(4\)}
\end{problem}

\begin{answer}
D
\end{answer}

\begin{solution}
由约束条件得 \(y\leqslant4-x\)，\(y\geqslant\dfrac{x+4}{3}\)，且 \(x\geqslant1\). 要使可行域非空，必须有
\[
\frac{x+4}{3}\leqslant4-x,
\]
即 \(1\leqslant x\leqslant2\). 可行域的顶点为 \(\left(1,3\right)\)，\(\left(1,\dfrac{5}{3}\right)\)，\((2,2)\). 在这三个点处，目标函数 \(z=3x-y\) 的值分别为 \(0\)，\(\dfrac{4}{3}\)，\(4\). 因此 \(z\) 的最大值为 \(4\)，故选 D.
\end{solution}

\begin{problem}
阅读程序框图，运行相应的程序，若输入 \(x\) 的值为 \(-4\)，则输出 \(y\) 的值为（\quad）

\bitmapfigure[max width=0.70\linewidth]{img/2011/tianjin_liberal/q3_fig1.png}

\choices
    {\(0.5\)}
    {\(1\)}
    {\(2\)}
    {\(4\)}
\end{problem}

\begin{answer}
C
\end{answer}

\begin{solution}
输入 \(x=-4\) 时，因 \(\lvert x\rvert=4>3\)，第一次循环后 \(x=\lvert-4-3\rvert=7\). 此时 \(\lvert x\rvert>3\)，第二次循环后 \(x=\lvert7-3\rvert=4\)；再次循环后 \(x=\lvert4-3\rvert=1\). 由于 \(\lvert1\rvert\leqslant3\)，循环结束，输出 \(y=2^1=2\). 故选 C.
\end{solution}

\begin{problem}
设集合 \(A = \{x \in \R \mid x - 2 > 0\}\)，\(B = \{x \in \R \mid x < 0\}\)，\(C = \{x \in \R \mid x(x - 2) > 0\}\)，则“\(x \in A \cup B\)”是“\(x \in C\)”的（\quad）

\choices
    {充分而不必要条件}
    {必要而不充分条件}
    {充分必要条件}
    {既不充分也不必要条件}
\end{problem}

\begin{answer}
C
\end{answer}

\begin{solution}
由定义，\(A=(2,+\infty)\)，\(B=(-\infty,0)\)，所以
\[
A\cup B=(-\infty,0)\cup(2,+\infty).
\]

又 \(x(x-2)>0\) 的解集为 \((-\infty,0)\cup(2,+\infty)\)，故
\[
x\in A\cup B\Longleftrightarrow x\in C.
\]

因此“\(x\in A\cup B\)”是“\(x\in C\)”的充分必要条件，故选 C.
\end{solution}

\begin{problem}
已知 \(a = \log_2 3.6\)，\(b = \log_4 3.2\)，\(c = \log_4 3.6\) 则（\quad）

\choices
    {\(a > b > c\)}
    {\(a > c > b\)}
    {\(b > a > c\)}
    {\(c > a > b\)}
\end{problem}

\begin{answer}
B
\end{answer}

\begin{solution}
由换底公式，\(c=\log_4 3.6=\dfrac{1}{2}\log_2 3.6=\dfrac{a}{2}\)，即 \(a=2c\). 又因 \(3.2<3.6\) 且底数 \(4>1\)，所以 \(b=\log_4 3.2<c\). 因此 \(a>c>b\)，故选 B.
\end{solution}

\begin{problem}
已知双曲线 \(\frac{x^2}{a^{2}} - \frac{y^2}{b^{2}} = 1 \) （\(a > 0, b > 0\)）的左顶点与抛物线 \(y^2 = 2px\) （\(p > 0\)） 的焦点的距离为4，且双曲线的一条渐近线与抛物线的准线的交点坐标为 \((-2, -1)\)，则双曲线的焦距为（\quad）

\choices
    {\(2\sqrt{3}\)}
    {\(2\sqrt{5}\)}
    {\(4\sqrt{3}\)}
    {\(4\sqrt{5}\)}
\end{problem}

\begin{answer}
B
\end{answer}

\begin{solution}
抛物线 \(y^2=2px\) 的焦点为 \(\left(\dfrac p2,0\right)\)，准线为 \(x=-\dfrac p2\). 由交点 \((-2,-1)\) 在准线上，得
\[
-\frac p2=-2,
\]
所以 \(p=4\). 双曲线左顶点为 \((-a,0)\)，由题意
\[
a+\frac p2=4,
\]
从而 \(a=2\). 双曲线的一条渐近线经过 \((-2,-1)\)，其斜率的绝对值为 \(\dfrac{1}{2}\)，故 \(\dfrac ba=\dfrac12\)，即 \(b=1\). 因此双曲线的半焦距为
\[
c=\sqrt{a^2+b^2}=\sqrt5,
\]
焦距为 \(2c=2\sqrt5\)，故选 B.
\end{solution}

\begin{problem}
已知函数 \(f(x) = 2\sin (\omega x + \varphi)\)，\(x \in \R\)，其中 \(\omega > 0\)，\(-\pi < \varphi \leqslant \pi\). 若 \( f(x) \) 的最小正周期为 \( 6\pi \)，且当 \( x = \frac{\pi}{2} \) 时，\( f(x) \) 取得最大值，则（\quad）

\choices
    {\(f(x)\) 在区间 \([-2\pi, 0]\) 上是增函数}
    {\(f(x)\) 在区间 \([-3\pi, -\pi]\) 上是增函数}
    {\(f(x)\) 在区间 \([3\pi, 5\pi]\) 上是减函数}
    {\(f(x)\) 在区间 \([4\pi, 6\pi]\) 上是减函数}
\end{problem}

\begin{answer}
A
\end{answer}

\begin{solution}
函数的最小正周期为 \(\dfrac{2\pi}{\omega}=6\pi\)，所以 \(\omega=\dfrac13\). 当 \(x=\dfrac\pi2\) 时函数取得最大值，故
\[
\sin\left(\frac\pi6+\varphi\right)=1.
\]
结合 \(-\pi<\varphi\leqslant\pi\)，得 \(\varphi=\dfrac\pi3\). 因而
\[
f(x)=2\sin\left(\frac x3+\frac\pi3\right).
\]
其导数的符号由 \(\cos\left(\dfrac x3+\dfrac\pi3\right)\) 决定，函数在
\[
\left(-\frac{5\pi}{2}+6k\pi,\frac\pi2+6k\pi\right)
\]
上递增，在 \(\left(\dfrac\pi2+6k\pi,\dfrac{7\pi}{2}+6k\pi\right)\) 上递减，其中 \(k\in\Z\). 区间 \([-2\pi,0]\) 包含于递增区间 \(\left(-\dfrac{5\pi}{2},\dfrac\pi2\right)\)，所以选 A.
\end{solution}

\begin{problem}
对实数 \(a\) 和 \(b\)，定义运算“\(\otimes\)”： \( a \otimes b = \begin{cases}
a, & a - b \leqslant 1,\\
b, & a - b > 1.
\end{cases} \) 设函数 \(f(x) = (x^{2} - 2) \otimes (x - 1)\)，\(x \in \R\). 若函数 \( y = f(x) - c \) 的图象与 \(x\) 轴恰有两个公共点，则实数 \(c\) 的取值范围是（\quad）

\choices
    {\((-1,1]\cup (2, + \infty)\)}
    {\((-2, - 1] \cup (1, 2]\)}
    {\((-\infty, -2) \cup (1, 2]\)}
    {\([-2, - 1]\)}
\end{problem}

\begin{answer}
B
\end{answer}

\begin{solution}
令 \(u=x^2-2\)，\(v=x-1\). 当 \(u-v\leqslant1\) 时
\[
x^2-x-2\leqslant0,
\]
即 \(-1\leqslant x\leqslant2\)，此时 \(f(x)=x^2-2\). 在其余区间 \(f(x)=x-1\). 因而
\[
f(x)=
\begin{cases}
x-1,&x<-1,\\
x^2-2,&-1\leqslant x\leqslant2,\\
x-1,&x>2.
\end{cases}
\]
方程 \(f(x)=c\) 的解的个数如下. 若 \(-2<c\leqslant-1\)，二次函数在 \([-1,2]\) 上给出两个解；若 \(1<c\leqslant2\)，二次函数给出一个解，且右侧一次函数给出一个解. 其余 \(c\) 值对应的解数均不是两个. 因此
\[
c\in(-2,-1]\cup(1,2].
\]
故选 B.
\end{solution}

\section{填空题}

\begin{problem}
已知集合 \(A = \{x \in \R \mid |x - 1| < 2\}\)，\(\Z\) 为整数集，则集合 \(A \cap \Z\) 中所有元素的和等于\fillinblank{}.
\end{problem}

\begin{answer}
3
\end{answer}

\begin{solution}
由 \(\lvert x-1\rvert<2\)，得 \(-1<x<3\). 其中的整数为 \(0\)，\(1\)，\(2\)，所以所有元素的和为
\[
0+1+2=3.
\]
\end{solution}

\begin{problem}
一个几何体的三视图如图所示（单位：m），则这个几何体的体积为\fillinblank{}\(\mathrm{m}^{3}\).

\bitmapfigure[max width=0.52\linewidth]{img/2011/tianjin_liberal/q10_fig1.png}
\end{problem}

\begin{answer}
4
\end{answer}

\begin{solution}
由三视图可知，该几何体由一个底面长为 \(2\)、宽为 \(1\)、高为 \(1\) 的长方体，以及其上方一个底面长为 \(1\)、宽为 \(1\)、高为 \(2\) 的长方体组成. 因此体积为
\[
2\times1\times1+1\times1\times2=4 \mathrm{~m}^{3}.
\]
\end{solution}

\begin{problem}
已知 \( \{a_{n}\} \) 为等差数列，\( S_{n} \) 为其前 \(n\) 项和. \( n \in \N^{*} \). 若 \( a_{3} = 16 \)，\( S_{20} = 20 \)，则 \(S_{10}\) 的值为\fillinblank{}.
\end{problem}

\begin{answer}
110
\end{answer}

\begin{solution}
设等差数列的首项为 \(a_1\)，公差为 \(d\). 由 \(a_3=16\)，得
\[
a_1+2d=16.
\]
又由 \(S_{20}=20\)，得
\[
\frac{20}{2}(2a_1+19d)=20,
\]
即 \(2a_1+19d=2\). 联立两式，得 \(d=-2\)，\(a_1=20\). 所以
\[
S_{10}=\frac{10}{2}\left(2a_1+9d\right)
=5(40-18)=110.
\]
\end{solution}

\begin{problem}
已知 \(\log_2 a + \log_2 b \geqslant 1\)，则 \(3^a + 9^b\) 的最小值为\fillinblank{}.
\end{problem}

\begin{answer}
18
\end{answer}

\begin{solution}
由对数有意义，\(a>0\)，\(b>0\). 由题设
\[
\log_2(ab)\geqslant1,
\]
所以 \(ab\geqslant2\). 令 \(X=a\ln3\)，\(Y=2b\ln3\)，则 \(X>0\)，\(Y>0\)，且
\[
XY=2ab(\ln3)^2\geqslant4(\ln3)^2.
\]
由基本不等式，
\[
3^a+9^b=\e^X+\e^Y
\geqslant2\e^{(X+Y)/2}
\geqslant2\e^{\sqrt{XY}}
\geqslant2\e^{2\ln3}=18.
\]
当且仅当 \(X=Y=2\ln3\)，即 \(a=2\)，\(b=1\) 时取等号. 因此最小值为 \(18\).
\end{solution}

\begin{problem}
如图，已知圆中两条弦 \(AB\) 与 \(CD\) 相交于点 \(F\)，\(E\) 是 \(AB\) 延长线上一点，且 \(DF = CF = \sqrt{2}\)，\(AF:FB:BE = 4:2:1\). 若 \(CE\) 与圆相切，则 \(CE\) 的长为\fillinblank{}.

\bitmapfigure[max width=0.42\linewidth]{img/2011/tianjin_liberal/q13_fig1.png}
\end{problem}

\begin{answer}
\(\frac{\sqrt{7}}{2}\)
\end{answer}

\begin{solution}
设 \(AF=4k\)、\(FB=2k\)、\(BE=k\)，其中 \(k>0\). 由相交弦定理，
\[
AF\cdot FB=CF\cdot DF=2,
\]
所以 \(8k^2=2\)，得 \(k=\dfrac12\). 因而 \(AF=2\)，\(FB=1\)，\(BE=\frac12\)，且 \(EA=AF+FB+BE=\frac72\).
由点 \(E\) 的切割线定理，\(CE\) 为切线，故
\[
CE^2=EA\cdot EB=\frac72\cdot\frac12=\frac74.
\]
因为 \(CE>0\)，所以 \(CE=\dfrac{\sqrt7}{2}\).
\end{solution}

\begin{problem}
已知直角梯形 \(ABCD\) 中，\(AD \parallel BC\)，\(\angle ADC = 90^{\circ}\)，\(AD = 2\)，\(BC = 1\)，\(P\) 是腰 \(DC\) 上的动点，则 \(\left|\overrightarrow{PA}+3\overrightarrow{PB}\right|\) 的最小值为\fillinblank{}.
\end{problem}

\begin{answer}
5
\end{answer}

\begin{solution}
建立直角坐标系，令 \(D=(0,0)\)，\(A=(0,2)\)，设 \(C=(t,0)\)，其中 \(t>0\). 由 \(AD\parallel BC\) 且 \(BC=1\)，可取 \(B=(t,1)\). 设 \(P=(s,0)\)，其中 \(0\leqslant s\leqslant t\). 则 \(\overrightarrow{PA}=(-s,2)\)，\(\overrightarrow{PB}=(t-s,1)\)，
从而
\[
\overrightarrow{PA}+3\overrightarrow{PB}=(3t-4s,5).
\]
因此
\[
\left|\overrightarrow{PA}+3\overrightarrow{PB}\right|
=\sqrt{(3t-4s)^2+25}\geqslant5.
\]
取 \(s=\dfrac{3t}{4}\)，此时 \(0\leqslant s\leqslant t\)，等号成立. 所以最小值为 \(5\).
\end{solution}

\section{解答题}

\begin{problem}
编号分别为 \(A_{1}\)，\(A_{2}\)，\(\cdots\)，\(A_{16}\) 名篮球运动员在某次训练比赛中的得分记录如下：

\begin{center}
\begin{tabular}{|c|c|c|c|c|c|c|c|c|}
\hline
运动员编号 & \(A_{1}\) & \(A_{2}\) & \(A_{3}\) & \(A_{4}\) & \(A_{5}\) & \(A_{6}\) & \(A_{7}\) & \(A_{8}\) \\
\hline
得分 & 15 & 35 & 21 & 28 & 25 & 36 & 18 & 34 \\
\hline
运动员编号 & \(A_{9}\) & \(A_{10}\) & \(A_{11}\) & \(A_{12}\) & \(A_{13}\) & \(A_{14}\) & \(A_{15}\) & \(A_{16}\) \\
\hline
得分 & 17 & 26 & 25 & 33 & 22 & 12 & 31 & 38 \\
\hline
\end{tabular}
\end{center}

\begin{enumerate}
    \item 将得分在对应区间内的人数填入相应的空格；
    \begin{center}
    \begin{tabular}{|c|c|c|c|}
    \hline
    区间 & \([10,20)\) & \([20,30)\) & \([30,40]\) \\
    \hline
    人数 &\fillinblank{}&\fillinblank{}&\fillinblank{} \\
    \hline
    \end{tabular}
    \end{center}
    \item 从得分在区间 \([20, 30)\) 内的运动员中随机抽取 2 人，
    \begin{enumerate}
    \item 用运动员编号列出所有可能的抽取结果；
    \item 求这 2 人得分之和大于 50 的概率.
    \end{enumerate}
\end{enumerate}
\end{problem}

\begin{solution}
\begin{enumerate}
\item 得分在 \([10,20)\) 内的运动员为 \(A_1\)，\(A_7\)，\(A_9\)，\(A_{14}\)，所以人数为 \(4\). 得分在 \([20,30)\) 内的运动员为 \(A_3\)，\(A_4\)，\(A_5\)，\(A_{10}\)，\(A_{11}\)，\(A_{13}\)，所以人数为 \(6\). 得分在 \([30,40]\) 内的运动员为 \(A_2\)，\(A_6\)，\(A_8\)，\(A_{12}\)，\(A_{15}\)，\(A_{16}\)，所以人数为 \(6\). 三个空格依次填 \(4\)，\(6\)，\(6\).

\item 得分在 \([20,30)\) 内的 6 名运动员为 \(A_3\)，\(A_4\)，\(A_5\)，\(A_{10}\)，\(A_{11}\)，\(A_{13}\)，任取 2 人的所有可能结果共有 \(\mathrm C_6^2=15\) 个，具体为
\[
\begin{gathered}
\{A_3,A_4\},\{A_3,A_5\},\{A_3,A_{10}\},\{A_3,A_{11}\},\{A_3,A_{13}\},\\
\{A_4,A_5\},\{A_4,A_{10}\},\{A_4,A_{11}\},\{A_4,A_{13}\},\\
\{A_5,A_{10}\},\{A_5,A_{11}\},\{A_5,A_{13}\},\\
\{A_{10},A_{11}\},\{A_{10},A_{13}\},\{A_{11},A_{13}\}.
\end{gathered}
\]

其中得分之和大于 \(50\) 的结果为
\[
\{A_4,A_5\},\ \{A_4,A_{10}\},\ \{A_4,A_{11}\},\ \{A_5,A_{10}\},\ \{A_{10},A_{11}\}.
\]
所以所求概率为
\[
\frac{5}{15}=\frac13.
\]
\end{enumerate}
\end{solution}

\begin{problem}
在 \(\triangle ABC\) 中，内角 \(A\)，\(B\)，\(C\) 的对边分别为 \(a\)，\(b\)，\(c\)，已知 \(B = C\)，\(2b = \sqrt{3}a\).
\begin{enumerate}
\item 求 \(\cos A\) 的值；
\item \(\cos \left(2A + \frac{\pi}{4}\right)\) 的值.
\end{enumerate}
\end{problem}

\begin{solution}
由 \(B=C\)，根据正弦定理可知 \(b=c\). 在三角形中使用余弦定理，得
\[
a^2=b^2+c^2-2bc\cos A=2b^2(1-\cos A).
\]
由 \(2b=\sqrt3a\)，得 \(b^2=\dfrac34a^2\). 代入上式并约去 \(a^2\)，得
\[
1=\frac32(1-\cos A),
\]
所以 \(\cos A=\dfrac13\). 因为 \(0<A<\pi\)，故
\[
\sin A=\sqrt{1-\cos^2A}=\frac{2\sqrt2}{3}.
\]
于是 \(\cos2A=2\cos^2A-1=-\frac79\)，
\(\sin2A=2\sin A\cos A=\frac{4\sqrt2}{9}\).
因此
\[
\cos\left(2A+\frac\pi4\right)
=\frac{\cos2A-\sin2A}{\sqrt2}
=-\frac{4}{9}-\frac{7\sqrt2}{18}
=-\frac{8+7\sqrt2}{18}.
\]
\end{solution}

\begin{problem}
如图，在四棱锥 \(P - ABCD\) 中，底面 \(ABCD\) 为平行四边形，\(\angle ADC = 45^{\circ}\)，\(AD = AC = 1\)，\(O\) 为 \(AC\) 中点，\(PO \perp\) 平面 \(ABCD\)，\(PO = 2\)，\(M\) 为 \(PD\) 中点.
\begin{enumerate}
\item 证明： \(PB \parallel\) 平面 \(ACM\)；
\item 证明： \(AD \perp\) 平面 \(PAC\)；
\item 求直线 \(AM\) 与平面 \(ABCD\) 所成角的正切值.
\end{enumerate}

\bitmapfigure[max width=0.42\linewidth]{img/2011/tianjin_liberal/q17_fig1.png}
\end{problem}

\begin{solution}
建立空间直角坐标系，令 \(D=(0,0,0)\)，\(A=(1,0,0)\). 因为 \(\angle ADC=45^\circ\)、\(AD=1\)、\(AC=1\)，可取 \(C=(1,1,0)\)，从而由平行四边形得 \(B=(2,1,0)\). 又 \(O=\left(1,\frac12,0\right)\)，\(P=\left(1,\frac12,2\right)\)，
\(M=\left(\frac12,\frac14,1\right)\).
\begin{enumerate}
\item 有 \(\overrightarrow{PB}=(1,\frac12,-2)\)，
\(\overrightarrow{AC}=(0,1,0)\)，
\(\overrightarrow{AM}=(-\frac12,\frac14,1)\).
直接计算得
\[
\overrightarrow{PB}=\overrightarrow{AC}-2\overrightarrow{AM}.
\]
所以 \(\overrightarrow{PB}\) 是平面 \(ACM\) 内两个相交方向向量的线性组合. 若 \(B\) 在平面 \(ACM\) 内，则存在实数 \(\lambda\)，\(\mu\)，使得
\[
\overrightarrow{AB}=\lambda\overrightarrow{AC}+\mu\overrightarrow{AM}.
\]
比较 \(z\) 坐标得 \(\mu=0\)，此时右端的 \(x\) 坐标为 \(0\)，不等于 \(\overrightarrow{AB}\) 的 \(x\) 坐标 \(1\)，矛盾. 因此 \(B\) 不在平面 \(ACM\) 内，故 \(PB\parallel\) 平面 \(ACM\).

\item 有 \(\overrightarrow{AD}=(-1,0,0)\)，
\(\overrightarrow{AC}=(0,1,0)\)，
\(\overrightarrow{AP}=(0,\frac12,2)\).
因此 \(\overrightarrow{AD}\cdot\overrightarrow{AC}=0\)，\(\overrightarrow{AD}\cdot\overrightarrow{AP}=0\). 直线 \(AD\) 垂直于平面 \(PAC\) 内相交的两条直线 \(AC\)、\(AP\)，所以 \(AD\perp\) 平面 \(PAC\).

\item 点 \(M\) 在平面 \(ABCD\) 上的射影为
\[
M'=\left(\frac12,\frac14,0\right).
\]
所以直线 \(AM\) 在底面内的射影为 \(AM'\)，且 \(AM'=\sqrt{\left(\frac12\right)^2+\left(\frac14\right)^2}=\frac{\sqrt5}{4}\)，
\(MM'=1\).
设直线 \(AM\) 与平面 \(ABCD\) 所成的角为 \(\theta\)，则
\[
\tan\theta=\frac{MM'}{AM'}=\frac{1}{\sqrt5/4}=\frac{4}{\sqrt5}.
\]
\end{enumerate}
\end{solution}

\begin{problem}
设椭圆 \(\frac{x^2}{a^{2}} + \frac{y^2}{b^{2}} = 1 \) （\(a > b > 0\)）的左、右焦点分别为 \(F_1\)，\(F_2\). 点 \(P(a, b)\) 满足 \(|PF_2| = |F_1F_2|\).
\begin{enumerate}
\item 求椭圆的离心率 \( e \)；
\item 设直线 \(PF_{2}\) 与椭圆相交于 \(A\)，\(B\) 两点，若直线 \(PF_{2}\) 与圆 \((x + 1)^{2} + (y - \sqrt{3})^{2} = 16\) 相交于 \(M\)，\(N\) 两点，且 \(|MN| = \frac{5}{8} |AB|\)，求椭圆的方程.
\end{enumerate}
\end{problem}

\begin{solution}
\begin{enumerate}
\item 设椭圆的半焦距为 \(c\)，则 \(c^2=a^2-b^2\)，且 \(F_1=(-c,0)\)、\(F_2=(c,0)\). 由 \(P=(a,b)\) 及题设，
\[
(a-c)^2+b^2=4c^2.
\]
利用 \(b^2=a^2-c^2\)，得
\[
2a^2-2ac=4c^2.
\]
令 \(e=\dfrac ca\)，并注意 \(e>0\)，则
\[
2-2e=4e^2,
\]
即 \(2e^2+e-1=0\). 所以 \(e=\dfrac12\).

\item 由 \(e=\dfrac12\)，得 \(c=\dfrac a2\)，\(b=\dfrac{\sqrt3}{2}a\). 直线 \(PF_2\) 的方向单位向量可取
\[
\bs{u}=\left(\frac12,\frac{\sqrt3}{2}\right).
\]
以 \(F_2\) 为参数原点，设直线上的点为
\[
(x,y)=F_2+t\bs{u}=\left(\frac a2+\frac t2,\frac{\sqrt3\,t}{2}\right),
\]
其中 \(t\) 表示有向距离. 代入椭圆方程，得
\[
\frac{(a+t)^2}{4a^2}+\frac{t^2}{a^2}=1,
\]
即
\[
5t^2+2at-3a^2=0.
\]
其两根为 \(t=-a\)、\(t=\dfrac{3a}{5}\)，所以
\[
|AB|=\frac{3a}{5}+a=\frac{8a}{5}.
\]

圆心为 \(Q=(-1,\sqrt3)\)，半径为 \(4\). 直线 \(PF_2\) 的方程为
\[
\sqrt3x-y-\frac{\sqrt3a}{2}=0.
\]
圆心到该直线的距离为
\[
d=\frac{\left|-\sqrt3-\sqrt3-\dfrac{\sqrt3a}{2}\right|}{2}
=\frac{\sqrt3(a+4)}{4}.
\]
所以圆在该直线上的弦长为
\[
|MN|=2\sqrt{16-d^2}
=2\sqrt{16-\frac{3(a+4)^2}{16}}.
\]
由 \(|MN|=\dfrac58|AB|=a\)，得
\[
a^2=64-\frac34(a+4)^2,
\]
即
\[
7a^2+24a-208=0.
\]
解得 \(a=4\) 或 \(a=-\dfrac{52}{7}\). 因 \(a>0\)，故 \(a=4\)，\(b=2\sqrt3\)，且此时 \(d=2\sqrt3<4\)，确有两个交点. 所求椭圆方程为
\[
\frac{x^2}{16}+\frac{y^2}{12}=1.
\]
\end{enumerate}
\end{solution}



\begin{problem}
已知函数 \(f(x) = 4x^{3} + 3tx^{2} - 6t^{2}x + t - 1\)，\(x \in \R\)，其中 \( t \in \R\).
\begin{enumerate}
\item 当 \(t = 1\) 时，求曲线 \(y = f(x)\) 在点 \((0, f(0))\) 处的切线方程；
\item 当 \(t \neq 0\) 时，求 \(f(x)\) 的单调区间；
\item 证明：对任意的 \(t \in (0, +\infty)\)，\(f(x)\) 在区间 \((0, 1)\) 内均存在零点.
\end{enumerate}
\end{problem}

\begin{solution}
\begin{enumerate}
\item 当 \(t=1\) 时，\(f(x)=4x^3+3x^2-6x\)，\(f(0)=0\).
又
\[
f'(x)=12x^2+6x-6,
\]
所以 \(f'(0)=-6\). 曲线在 \((0,f(0))=(0,0)\) 处的切线方程为
\[
y=-6x.
\]

\item 对任意 \(t\)，
\[
f'(x)=12x^2+6tx-6t^2
=6(x+t)(2x-t).
\]
当 \(t>0\) 时，\(-t<\dfrac t2\)，所以 \(f'(x)\) 在 \((-\infty,-t)\)、\(\left(\dfrac t2,+\infty\right)\) 上为正，在 \(\left(-t,\dfrac t2\right)\) 上为负. 因此 \(f(x)\) 的增区间为 \((-\infty,-t)\)、\(\left(\dfrac t2,+\infty\right)\)，减区间为 \(\left(-t,\dfrac t2\right)\).

当 \(t<0\) 时，\(\dfrac t2<-t\)，所以 \(f'(x)\) 在 \(\left(-\infty,\dfrac t2\right)\)、\((-t,+\infty)\) 上为正，在 \(\left(\dfrac t2,-t\right)\) 上为负. 因此 \(f(x)\) 的增区间为 \(\left(-\infty,\dfrac t2\right)\)、\((-t,+\infty)\)，减区间为 \(\left(\dfrac t2,-t\right)\).

\item 先计算
\[
f\left(\frac12\right)=-3t^2+\frac74t-\frac12
= -3\left(t-\frac{7}{24}\right)^2-\frac{47}{192}<0.
\]
当 \(0<t\leqslant1\) 时，由 \(t^2\leqslant t\)，有
\[
f(1)=3+4t-6t^2\geqslant3-2t\geqslant1>0.
\]
于是由零点存在定理，\(f(x)\) 在 \(\left(\dfrac12,1\right)\) 内存在零点. 当 \(t>1\) 时，\(f(0)=t-1>0\)，结合 \(f\left(\dfrac12\right)<0\)，由零点存在定理，\(f(x)\) 在 \(\left(0,\dfrac12\right)\) 内存在零点. 所以对任意 \(t>0\)，\(f(x)\) 在 \((0,1)\) 内均存在零点.
\end{enumerate}
\end{solution}

\begin{problem}
已知数列 \(\{a_{n}\}\) 与 \(\{b_{n}\}\) 满足： \(b_{n+1}a_{n} + b_{n}a_{n+1} = (-2)^{n} + 1\)，\(b_{n} = \frac{3 + (-1)^{n-1}}{2}\)，\(n \in \N^{*}\)，且 \(a_{1} = 2\).
\begin{enumerate}
\item 求 \(a_{2}\)，\(a_{3}\) 的值；
\item 设 \(c_{n} = a_{2n + 1} - a_{2n - 1}\)，\(n \in \N^{*}\)，证明：数列 \(\{c_{n}\}\) 是等比数列；
\item 设 \(S_{n}\) 为 \(\{a_{n}\}\) 的前 \(n\) 项和，证明 \(\frac{S_{1}}{a_{1}} + \frac{S_{2}}{a_{2}} + \ldots + \frac{S_{2n-1}}{a_{2n-1}} + \frac{S_{2n}}{a_{2n}} \leqslant n - \frac{1}{3}\) （\(n \in \N^{*}\)）.
\end{enumerate}
\end{problem}

\begin{solution}
\begin{enumerate}
\item 由 \(b_1=2\)，\(b_2=1\)，\(b_3=2\).
在原递推式中令 \(n=1\)，得
\[
a_1+2a_2=-1.
\]
代入 \(a_1=2\)，得 \(a_2=-\dfrac32\). 令 \(n=2\)，得
\[
2a_2+a_3=5,
\]
所以 \(a_3=8\).

\item 对 \(n\in\N^*\)，分别在原递推式中取 \(2n-1\) 和 \(2n\)，得
\[
a_{2n-1}+2a_{2n}=1-2\cdot4^{n-1},
\]
\[
2a_{2n}+a_{2n+1}=4^n+1.
\]
两式相减，得
\[
c_n=a_{2n+1}-a_{2n-1}
=4^n+2\cdot4^{n-1}
=6\cdot4^{n-1}.
\]
因此
\[
\frac{c_{n+1}}{c_n}=4,
\]
所以数列 \(\{c_n\}\) 是以 \(4\) 为公比的等比数列.

\item 由上一步的结果及 \(a_1=2\)，得
\[
a_{2n-1}=a_1+\sum_{j=1}^{n-1}c_j
=2+\sum_{j=1}^{n-1}6\cdot4^{j-1}
=2\cdot4^{n-1}.
\]
再由 \(a_{2n-1}+2a_{2n}=1-2\cdot4^{n-1}\)，得
\[
a_{2n}=\frac{1-4^n}{2}.
\]
于是
\[
S_{2n-1}=\sum_{j=1}^{n}a_{2j-1}+\sum_{j=1}^{n-1}a_{2j}
=2\cdot4^{n-1}+\frac{n-1}{2},
\]
\[
S_{2n}=S_{2n-1}+a_{2n}=\frac n2.
\]
从而 \(\frac{S_{2n-1}}{a_{2n-1}}=1+\frac{n-1}{4^n}\)，
\(\frac{S_{2n}}{a_{2n}}=\frac{n}{1-4^n}\).
令
\[
q_n=\frac{S_{2n-1}}{a_{2n-1}}+\frac{S_{2n}}{a_{2n}}
=1+\frac{n-1}{4^n}-\frac{n}{4^n-1}.
\]
因为
\[
\frac{n-1}{4^n}<\frac{n}{4^n-1},
\]
所以 \(q_n<1\). 又 \(q_1=1-\dfrac13=\dfrac23\). 因此
\[
\sum_{k=1}^{2n}\frac{S_k}{a_k}
=\sum_{j=1}^{n}q_j
\leqslant\frac23+(n-1)
=n-\frac13.
\]
\end{enumerate}
\end{solution}
